\documentclass[11pt,a4paper]{article}

\usepackage[margin=1in]{geometry}
\usepackage{amsmath,amssymb,amsfonts}
\usepackage{mathtools}
\usepackage{lmodern}
\usepackage[T1]{fontenc}
\usepackage[utf8]{inputenc}
\usepackage{textcomp}
\usepackage{microtype}
\usepackage{parskip}
\usepackage{enumitem}
\usepackage{array}
\usepackage{booktabs}
\usepackage{hyperref}
\hypersetup{hidelinks,
  pdftitle={From Primitives to the Rate of Becoming},
  pdfauthor={Allen Proxmire}}

\setlength{\parindent}{0pt}
\setlength{\parskip}{6pt plus 2pt minus 1pt}

\title{\vspace{-1cm}\Large\bfseries From Primitives to the Rate of Becoming\\[6pt]
\large\mdseries Slow Light, Stopped Light, Optical Lattice Clocks, and Gravitational Redshift from Substrate Primitives}
\author{Allen Proxmire}
\date{May 2026}

\begin{document}
\maketitle
\thispagestyle{empty}

\section{The Question}

\subsection*{What this walkthrough derives}

Three experimental programs sit at the frontier of atomic, molecular, and optical (AMO) physics:

\begin{enumerate}[noitemsep]
  \item \textbf{Hau (1999)}: light propagating through an ultracold atomic cloud under electromagnetically-induced-transparency (EIT) conditions has its group velocity reduced from $c \approx 3\times 10^8$ m/s to $\sim 17$ m/s --- a reduction by factor $\sim 10^7$.
  \item \textbf{Hau (2001)}: in the same EIT system, adiabatic ramp of the control field to zero stores the optical pulse coherently in atomic ground-state coherence; subsequent reactivation retrieves the pulse with phase information preserved. Light is stopped, then released.
  \item \textbf{Katori (2003) and Ye (et seq.)}: optical lattice clocks based on the ${}^1S_0 \leftrightarrow {}^3P_0$ transition in alkaline-earth atoms (Sr, Yb) reach fractional frequency stabilities of $\sim 10^{-19}$, well past cesium-fountain microwave clocks. Networked clocks at millimeter-scale separation resolve gravitational redshift directly.
\end{enumerate}

The walkthrough derives, from substrate primitives:

\begin{itemize}[noitemsep]
  \item The \textbf{substrate-level definition of local rate of becoming} as rate of phase accumulation in pre-individuation amplitudes.
  \item The \textbf{speed of light} $c$ as the coarse-grained expression of this rate; substrate $c$ is a fundamental constant.
  \item \textbf{Slow light} as group-velocity reduction of a dressed polariton, with substrate $c$ unchanged.
  \item \textbf{Stopped light} as adiabatic identity transfer between channel types.
  \item \textbf{Optical lattice clocks} as direct probes of local becoming rate via phase accumulation.
  \item \textbf{Gravitational redshift} as ED-gradient curvature in the becoming rate.
  \item \textbf{Clock networks} as differential ED-gradient sensors.
\end{itemize}

The walkthrough is fully self-contained: every required EIT algebra, dark-state-polariton derivation, adiabatic-rotation step, magic-wavelength condition, and gravitational-redshift identity is derived inline.

\subsection*{What Event Density claims}

All three results --- slow light, stopped light, optical lattice clocks --- are FORCED at the substrate level by a single foundational primitive: \textbf{the local rate of becoming}, defined precisely as the rate of phase accumulation in pre-individuation amplitudes. The mechanism is:

\begin{itemize}[noitemsep]
  \item Substrate $c$ is the fundamental rate at which the rule-type structure propagates. It is a constant.
  \item A medium dresses the photon's participation rule with atomic-coherence rule structure, producing a polariton --- a mixed-channel-type rule structure whose group velocity is the coarse-grained becoming rate of the dressed object, generally less than $c$.
  \item Adiabatic switching of the medium's coupling rotates the polariton's identity continuously between photonic and atomic-spin-coherence components.
  \item Atomic clock transitions are direct probes of the local becoming rate.
  \item Gravitational redshift is the ED-gradient curvature in the local becoming rate.
  \item Clock networks measure differential ED-gradient curvature with millimeter-scale resolution.
\end{itemize}

The substrate-level mechanism, in one sentence: \textbf{all three Hau--Katori--Ye phenomena manipulate or measure the local rate of becoming, which is the substrate-level rate of phase accumulation that coarse-grains to both the speed of light and the tick rate of atomic clocks}.

\section{The Primitives}

\subsection{P-RB-1. Local rate of becoming}

The \emph{local rate of becoming} at a substrate-level point $x$ is the rate at which a chain's pre-individuation amplitude at $x$ accumulates phase per unit substrate-tick. Operationally, for a chain whose pre-individuation amplitude is $\psi(x, t) = \sum_n c_n(t) e^{-i\theta_n(t)}|r_n(x)\rangle$ aligned with eigen-rule $|r_n(x)\rangle$ of eigen-energy $E_n(x)$, the becoming rate is
\[
\omega_{\mathrm{becoming}}(x) \equiv d\theta_n/dt = E_n(x)/\hbar.
\]

\textbf{Substrate-level meaning.} The pre-individuation amplitude is the substrate-level superposition of consistent rule continuations weighted by complex amplitudes. Phase accumulation in this amplitude is the substrate-level analog of ``passage of time'' for the chain's identity: each phase increment corresponds to one substrate-tick of becoming.

\textbf{Two coarse-grained quantities.} When coarse-grained to standard observables:
\begin{itemize}[noitemsep]
  \item For massless chains (light), the rate of becoming sets the propagation speed: $c$ is the coarse-grained becoming rate of the substrate's rule-type structure propagating freely.
  \item For massive chains, the rate of becoming is the standard QM angular frequency $\omega = E/\hbar$.
\end{itemize}

\textbf{This is the load-bearing primitive of the walkthrough.}

\subsection{P-RB-2. Light-like chain}

A \emph{light-like chain} is a minimal ED-channel (multiplicity $M = 1$) whose participation rule has zero rest energy: $E = \hbar\omega = \hbar c|\mathbf{k}|$.

\textbf{Substrate $c$ is a constant.} This is non-negotiable. Slow light, refractive index, gravitational time dilation, and other apparent modifications of light's propagation rate are all \emph{coarse-grained dressed-object} effects. The substrate-level photon always propagates at $c$.

\subsection{P-RB-3. ED-gradient}

An \emph{ED-gradient} is a spatial variation in the local rate of becoming: $\nabla_x\omega_{\mathrm{becoming}}(x) \neq 0$.

Sources include coupling to a medium, gravitational potential gradients, and spatially-varying field configurations. Higher-order spatial variation produces ED-gradient curvature.

\subsection{P-RB-4. Identity alignment}

The \emph{identity alignment} of a chain at substrate-level position $x$ is its commitment to one specific participation rule from the local alignment set $\mathcal{R}(x) = \{r_n(x)\}$. Smooth changes in the substrate's rule structure can preserve a chain's identity alignment if the change is slow enough (adiabatic).

\subsection{P-RB-5. Medium-induced rule coupling}

A \emph{medium} is a substrate region containing additional chains whose rule-type structure couples to a propagating chain's rule structure. When a photonic chain $\hat a$ at frequency $\omega$ couples to atomic chains in a 3-level $\Lambda$ configuration with control field $\hat\Omega_c$, the effective interaction is
\[
H_{\mathrm{coupling}} = g(\hat a^\dagger\sigma_{eg} + \hat a\sigma_{ge}) + \Omega_c(\sigma_{em} + \sigma_{me}) + \mathrm{h.c.}
\]
This produces \emph{polaritons} --- mixed-channel-type rule structures, derived inline in \S4.

\subsection{P-RB-6. Coherence storage}

\emph{Coherence storage} is the substrate-level transfer of a chain's identity alignment from one channel type (photonic) to another (atomic spin coherence) via adiabatic rotation of the participation-rule structure. The adiabatic condition: $|d/dt(\Omega_c, g)| \ll \Delta_{\mathrm{gap}}$.

\subsection{P-RB-7. Local becoming measurement}

\emph{Local becoming measurement} is any operational protocol that probes the local rate of becoming directly. The atomic-clock transition is the canonical example: a chain in superposition between two long-lived eigen-rules accumulates relative phase at rate $(E_2 - E_1)/\hbar$, which is the local becoming-rate difference between the two rules.

\textbf{No additional substrate primitives are introduced beyond P-RB-1 through P-RB-7.}

\section{Deriving the Speed of Light from the Rate of Becoming}

\subsection{Substrate $c$ as V1-kernel propagation rate}

The substrate's V1 kernel encodes how rule-type structure propagates between adjacent substrate points per substrate-tick. The propagation rate is a substrate-level constant: in the coarse-grained limit, the V1 kernel reduces to a $\delta(x - ct)$ propagation pattern.

Standard physics identifies this propagation rate with the speed of light $c \approx 3\times 10^8$ m/s. The substrate-level statement is: $c$ is the V1-kernel propagation rate, and it is invariant under all coordinate, gauge, and material transformations.

\subsection{Group velocity vs phase velocity}

For a chain propagating freely in vacuum, both phase velocity ($v_\phi = \omega/|\mathbf{k}|$) and group velocity ($v_g = d\omega/d|\mathbf{k}|$) equal $c$. For a chain propagating in a medium, the dispersion relation is modified. The chain's coarse-grained becoming rate becomes a non-linear function $\omega(|\mathbf{k}|)$, and the group velocity $v_g = d\omega/d|\mathbf{k}|$ can be much less than $c$. \textbf{The substrate $c$ is unchanged}; the modification lives entirely in the dispersion relation of the dressed object.

\subsection{Refractive index from ED-gradient slope}

The refractive index $n$ is defined operationally by $n(\omega) = c/v_{\mathrm{phase}}(\omega)$. In the substrate-level picture, $n$ is the coarse-grained image of an ED-gradient slope.

\textbf{Substrate-level statement.} Refractive index is not a property of an underlying medium that ``slows'' light. It is the coarse-grained image of how the dressed-object's rule structure couples the photonic channel to the medium's atomic-coherence channel.

\subsection{Why substrate $c$ must be invariant}

If $c$ varied with environment, the substrate's V1-kernel propagation would not be a single rule-type primitive --- it would be a family of rules. This would violate the substrate's primitive economy: the kernel would no longer be one primitive.

\section{Slow Light as Mixed-Channel Rule Structure (Hau 1999)}

\subsection{Setup: the 3-level $\Lambda$ system}

Hau's slow-light experiment uses ultracold Na atoms with three relevant atomic levels. A weak probe field at frequency $\omega_p$ couples $|g\rangle \leftrightarrow |e\rangle$. A strong control field at frequency $\omega_c$ couples $|m\rangle \leftrightarrow |e\rangle$. Two-photon resonance: $\omega_p - \omega_c$ matches the $|g\rangle \to |m\rangle$ ground-state splitting.

\subsection{Three-level Hamiltonian}

In the rotating-wave approximation:
\[
H_{\mathrm{atom}} = -\Delta_p|e\rangle\langle e| + g\hat a|e\rangle\langle g| + \Omega_c|e\rangle\langle m| + \mathrm{h.c.}
\]
For $N$ atoms, the collective coupling to the photon mode is $g\sqrt N$.

\subsection{Dark state and EIT}

The atomic Hamiltonian admits a \emph{dark state} --- an eigenstate that does not couple to the excited level:
\[
|D\rangle = \cos\theta|g\rangle - \sin\theta|m\rangle, \quad \tan\theta = g\sqrt N\hat a/\Omega_c.
\]
Verification: the matrix element of $H_{\mathrm{atom}}$ between $|D\rangle$ and $|e\rangle$ is
\[
\langle e|H_{\mathrm{atom}}|D\rangle = g\sqrt N\hat a\cos\theta - \Omega_c\sin\theta = 0
\]
for the given $\tan\theta$. The dark state is transparent to the resonant absorption process. This is the EIT condition.

\subsection{Dark-state polariton (mixed-channel rule structure)}

The full system admits an eigenstate that is a coherent superposition of photonic and atomic-coherence components:
\[
|\Psi_{\mathrm{polariton}}\rangle = \cos\theta|1_{\mathrm{photon}}, g\rangle - \sin\theta|0_{\mathrm{photon}}, m\text{-coherence}\rangle,
\]
with mixing angle $\tan\theta = g\sqrt N/\Omega_c$.

\textbf{Substrate-level meaning.} The polariton is one substrate object whose rule structure is composite. It is not a photon traveling alongside an atomic excitation; it is one chain whose identity alignment spans two channel types.

\subsection{Group velocity of the polariton}

The polariton's effective Hamiltonian at two-photon resonance is
\[
H_{\mathrm{polariton}} = c k\cos^2\theta.
\]
The group velocity is
\[
v_g = dE_{\mathrm{polariton}}/d(\hbar k) = c\cos^2\theta = c\cdot\frac{\Omega_c^2}{\Omega_c^2 + g^2 N}.
\]
For weak control field ($\Omega_c \ll g\sqrt N$): $v_g \to c\cdot\Omega_c^2/(g^2 N) \ll c$.

For typical Hau-experiment values ($g\sqrt N \sim 10$ GHz, $\Omega_c \sim 10$ MHz): $v_g \sim 10$--30 m/s. \textbf{Hau measured $\sim 17$ m/s.}

\subsection{Substrate-level reading}

\begin{itemize}[noitemsep]
  \item The propagating object is a polariton, not a substrate-level photon.
  \item The polariton's group velocity is the coarse-grained becoming rate of the dressed mode.
  \item When the photonic component dominates ($\cos^2\theta \approx 1$), the polariton becomes light-like.
  \item When the atomic component dominates ($\sin^2\theta \approx 1$), the polariton becomes essentially stationary.
  \item \textbf{Substrate $c$ is unchanged.}
\end{itemize}

\section{Stopped Light as Identity Transfer (Hau 2001)}

\subsection{Adiabatic ramp of the control field}

Slowly reduce the control-field Rabi frequency: $\Omega_c(t) \to 0$ over a timescale long compared to the gap between the dark-state polariton and the excited atomic states. The mixing angle evolves: $\tan\theta(t) = g\sqrt N/\Omega_c(t) \to \infty$, so $\theta(t) \to \pi/2$. The polariton rotates from photonic-dominated to atomic-dominated:
\[
|\Psi_{\mathrm{polariton}}\rangle \to -|0_{\mathrm{photon}}, m\text{-coherence}\rangle \text{ as } \Omega_c \to 0.
\]

\subsection{Adiabatic identity preservation}

Provided the ramp is slow enough, the chain's identity alignment is preserved. \textbf{The chain's identity is now stored entirely in the medium's rule structure.} No photon is propagating; the chain's identity exists as ground-state atomic coherence.

\subsection{Storage time and decoherence}

The storage time is limited by the lifetime of the atomic ground-state coherence --- typically much longer than the excited-state lifetime. In Hau's experiment, storage times of $\sim 1$ ms were demonstrated.

\subsection{Retrieval via reversed ramp}

Re-applying the control field reverses the polariton rotation. The chain's identity returns to a photonic-dominated polariton, which propagates out of the medium with the original frequency, polarization, and phase information preserved.

\textbf{Substrate-level reading.} Stopped light is adiabatic identity transfer between channel types. No state is ``frozen''; no photon is ``captured.'' The substrate-level mechanism is continuous identity transfer, parameterized by the control-field strength.

\section{Optical Lattice Clocks as Becoming-Rate Probes (Katori, Ye)}

\subsection{The clock transition}

Strontium-87 has a metastable ${}^3P_0$ state above the ${}^1S_0$ ground state. The transition ${}^1S_0 \leftrightarrow {}^3P_0$ at 698 nm has natural linewidth $\sim 1$ mHz.

The chain's pre-individuation amplitude in the clock's superposition is
\[
\psi_{\mathrm{clock}}(t) = \frac{1}{\sqrt 2}\bigl[|{}^1S_0\rangle + e^{-i\omega_{\mathrm{clock}} t}|{}^3P_0\rangle\bigr],
\]
where $\omega_{\mathrm{clock}} = (E_{{}^3P_0} - E_{{}^1S_0})/\hbar \approx 2\pi\cdot 4.3\times 10^{14}$ Hz.

\subsection{Clock observable as becoming-rate measurement}

The clock's operational observable is the rate at which the relative phase between the two state components accumulates:
\[
d\theta_{\mathrm{relative}}/dt = \omega_{\mathrm{clock}} = (E_{{}^3P_0} - E_{{}^1S_0})/\hbar.
\]
By P-RB-1, the clock observable \emph{is} the local becoming rate of the clock-state participation rule, made operationally accessible via interferometry between the two state components.

\textbf{The clock measures local rate of becoming directly.}

\subsection{Magic wavelength as substrate-level Stark cancellation}

When the clock atoms are trapped in an optical lattice, the lattice light induces AC Stark shifts on both clock states:
\[
\Delta\omega_{\mathrm{clock}}(\lambda) = -\frac{I}{2\hbar}\bigl[\alpha_{{}^3P_0}(\lambda) - \alpha_{{}^1S_0}(\lambda)\bigr].
\]
At the \emph{magic wavelength} $\lambda_m \approx 813$ nm for Sr, $\alpha_{{}^3P_0}(\lambda_m) = \alpha_{{}^1S_0}(\lambda_m)$, so $\Delta\omega_{\mathrm{clock}}(\lambda_m) = 0$.

\textbf{Substrate-level statement.} At $\lambda_m$, the substrate-level photon-atom rule-type coupling produces \emph{no net rule-type-energy shift} on the clock-transition rule structure.

\subsection{Lamb-Dicke regime as motional-sideband suppression}

When the trapping potential is strong enough that $\eta \equiv k_{\mathrm{clock}}\cdot a_{\mathrm{vib}} \ll 1$ (Lamb-Dicke parameter), the atomic motion does not modulate the clock-transition phase accumulation. Doppler shifts and motional sidebands are suppressed.

\subsection{Clock precision as substrate-level isolation}

The optical lattice clock's fractional frequency stability is
\[
\sigma_y(\tau) \sim 1/(\omega_{\mathrm{clock}}\sqrt{N\cdot T_{\mathrm{coh}}\cdot\tau}).
\]
The current frontier is $\sigma_y \sim 10^{-19}$ over hour-scale integration.

\subsection{Why clocks measure geometry of becoming, not ``time''}

In the standard reading, an atomic clock measures time. In the substrate-level reading, the clock measures \emph{the local rate of becoming} --- the rate of phase accumulation in the chain's pre-individuation amplitude at the clock's spatial position. \textbf{The substrate-level claim:} what optical lattice clocks measure with $10^{-19}$ precision is not a universal ``time'' but a \emph{local geometric quantity}.

\section{Gravitational Redshift as ED-Gradient Curvature}

\subsection{Gravitational potential as ED-gradient}

In the weak-field limit, the gravitational potential $\Phi(\mathbf{x})$ is the coarse-grained image of an ED-gradient in the local rate of becoming:
\[
\omega_{\mathrm{becoming}}(\mathbf{x}) = \omega_0[1 + \Phi(\mathbf{x})/c^2],
\]
where $\omega_0$ is the becoming rate at infinite distance from any mass.

\subsection{Local becoming rate $d\tau$}

A clock at position $\mathbf{x}$ measures
\[
\omega_{\mathrm{becoming}}(\mathbf{x})\cdot dt = (E/\hbar)(1 + \Phi/c^2)\, dt = (E/\hbar)\, d\tau,
\]
where $d\tau = (1 + \Phi/c^2)\, dt$ is the \emph{proper-time interval}. In the weak-field limit:
\[
d\tau \approx \sqrt{1 + 2\Phi/c^2}\, dt = \sqrt{g_{00}}\, dt,
\]
matching the standard general-relativistic identity.

\subsection{Redshift between two heights}

Two clocks at heights $h_1$ and $h_2$ have local potentials $\Phi_i = g h_i$. Their becoming-rate ratio is
\[
\omega_1/\omega_2 \approx 1 + g(h_1 - h_2)/c^2.
\]
For $h_1 > h_2$: $\omega_1 > \omega_2$, so the upper clock ticks faster:
\[
\Delta\omega/\omega = -g\Delta h/c^2 \quad \text{(gravitational redshift formula)}.
\]

\subsection{Substrate-level reading}

Gravitational redshift is not a property of light; it is a property of the substrate's becoming-rate gradient. The redshift is FORCED by the ED-gradient in $\omega_{\mathrm{becoming}}(\mathbf{x})$.

\textbf{Optical lattice clocks detect ED-curvature directly.} For $\Delta h = 1$ mm and Earth's $g \approx 10$ m/s$^2$: $\Delta\omega/\omega \approx 10^{-19}$. This is the precision frontier of current clocks.

\subsection{Why GR identifies $g_{00}$ with becoming-rate ratio}

The ED-substrate-gravity reading: $g_{00}$ is the coarse-grained image of $\omega_{\mathrm{becoming}}/\omega_0$, the ratio of local becoming rate to the reference becoming rate at infinite distance.

\section{Clock Networks as ED-Gradient Sensors}

\subsection{Differential clock measurements}

A network of optical lattice clocks at different spatial positions can be compared via fiber-optic frequency dissemination. The differential measurement
\[
\Delta\nu_{ij} = \nu_i - \nu_j
\]
cancels common-mode noise and isolates the \emph{position-dependent} becoming-rate difference.

\subsection{Sensitivity to ED-gradients}

For sufficiently isolated clocks, the differential signal becomes dominated by the gravitational potential difference $\Delta\Phi/c^2$. Current state-of-the-art Sr clock networks can resolve potential differences corresponding to $\sim 1$ mm altitude differences over hour-scale integration.

\subsection{Substrate-level mapping}

Each clock measures its local becoming rate. A network of clocks samples $\omega_{\mathrm{becoming}}$ at many points; differential measurements reveal the spatial structure --- the substrate's ED-gradient geometry. Clock networks are \emph{ED-geometry sensors}.

\subsection{Why this is more than gravimetry}

Standard gravimeters measure $g(\mathbf{x})$. Clock networks measure gravitational \emph{potential} $\Phi(\mathbf{x})$ directly via the becoming-rate gradient. Clock networks have access to the more fundamental substrate-level field.

\subsection{Future tests}

\begin{itemize}[noitemsep]
  \item \textbf{Variations of fundamental constants.}
  \item \textbf{Dark-matter searches} via oscillations in fundamental constants.
  \item \textbf{Gravitational-wave detection} as transient signals in differential measurements.
\end{itemize}

\section{Substrate-Level Reading}

\begin{center}
\footnotesize
\begin{tabular}{@{}p{5.5cm}p{9cm}@{}}
\toprule
Standard AMO/GR object & Substrate-level meaning \\
\midrule
Speed of light $c$ & Coarse-grained becoming rate of free-photon rule structure \\
Refractive index $n(\omega)$ & ED-gradient slope of dressed becoming-rate function \\
Group velocity $v_g$ & Coarse-grained becoming rate of the dressed (polariton) chain \\
Slow light & Polariton becoming-rate reduction via mixing with atomic-coherence \\
Stopped light & Identity adiabatically transferred from photonic to atomic-coherence channel \\
Atomic transition frequency & Local becoming rate; coarse-grains to $E/\hbar$ \\
Magic wavelength & Substrate-level Stark cancellation on clock-transition rule structure \\
Lamb-Dicke regime & Decoupling of internal-state from external-motion rule structure \\
Optical lattice clock & Direct probe of local rate of becoming \\
Gravitational potential $\Phi(\mathbf{x})$ & Coarse-grained ED-gradient in local becoming rate \\
Gravitational redshift & Ratio of becoming rates at two points \\
Time dilation $d\tau = \sqrt{g_{00}}\, dt$ & Local becoming-rate ratio relative to reference \\
Clock network & Spatially-distributed sensor of ED-geometry \\
\bottomrule
\end{tabular}
\end{center}

\subsection{The unified phenomenon}

The Hau, Katori, and Ye experimental programs are a single substrate-level enterprise: they manipulate or measure the local rate of becoming. The substrate-level mechanism is the same in all three: the local rate of becoming is the foundational substrate primitive; light propagation, atomic clock ticking, and gravitational time dilation are different coarse-grained expressions of the same underlying becoming-rate field.

\section{What's Forced, What's Inherited, What's Open}

\subsection*{What is FORCED by the substrate ontology}

\begin{itemize}[noitemsep]
  \item \textbf{Local rate of becoming as rate of phase accumulation in pre-individuation amplitudes} (P-RB-1).
  \item \textbf{Speed of light $c$ as constant V1-kernel propagation rate}.
  \item \textbf{Group velocity reduction in slow light} via mixed-channel rule structure.
  \item \textbf{Stopped light as adiabatic identity transfer}.
  \item \textbf{Optical-clock observable as local rate of becoming} via P-RB-1.
  \item \textbf{Magic-wavelength condition} as substrate-level Stark cancellation.
  \item \textbf{Gravitational redshift as ED-gradient curvature}.
  \item \textbf{Identity of GR's $g_{00}$ with becoming-rate ratio}.
\end{itemize}

\subsection*{What is FORM-FORCED-INHERITED (and re-derived inside this document)}

Three-level $\Lambda$-system Hamiltonian in rotating-wave approximation; dark-state structure in EIT; dark-state polariton composition; group velocity formula; AC Stark shift formula; Lamb-Dicke parameter; adiabatic-evolution identity preservation; weak-field metric expansion; Lindblad master-equation framework.

\subsection*{What remains OPEN}

\begin{itemize}[noitemsep]
  \item Many-body ED-flow beyond mean-field treatment.
  \item Nonlinear becoming gradients (hyperpolarizability shifts).
  \item Relativistic corrections to becoming rate in strong-field regimes.
  \item Quantum-coherence-enhanced clocks reaching Heisenberg-limited precision.
  \item Frequency-comb-based clock comparisons.
  \item Drift of fundamental constants.
  \item Tests of Lorentz/CPT invariance.
\end{itemize}

\section{What This Argument Establishes}

\textbf{Claim 1.} The local rate of becoming, defined as the rate of phase accumulation in pre-individuation amplitudes, is a foundational substrate primitive (P-RB-1). It coarse-grains to standard QM angular frequency $E/\hbar$ for atomic-state superpositions and to substrate $c$ for free-photon propagation.

\textbf{Claim 2.} The substrate-level $c$ is a constant --- the V1-kernel propagation rate. Slow light, refractive index, and other apparent modifications are coarse-grained dressed-object effects; substrate $c$ is invariant.

\textbf{Claim 3.} Slow light (Hau 1999) is FORCED as group-velocity reduction of the dark-state polariton, with photon-component weight $\cos^2\theta = \Omega_c^2/(\Omega_c^2 + g^2 N)$. Group velocity $v_g = c\cos^2\theta$.

\textbf{Claim 4.} Stopped light (Hau 2001) is FORCED as adiabatic identity transfer from photonic to atomic-coherence channel as $\Omega_c \to 0$. Retrieval is the reverse rotation; no information is lost.

\textbf{Claim 5.} Optical lattice clocks (Katori, Ye) measure the local rate of becoming directly via interferometry. Magic wavelength is the substrate-level condition for zero net rule-type-energy shift; Lamb-Dicke regime decouples internal-state from motional rule structure.

\textbf{Claim 6.} Gravitational redshift is FORCED by the ED-gradient in local becoming rate: $\omega_{\mathrm{becoming}}(\mathbf{x}) = \omega_0(1 + \Phi(\mathbf{x})/c^2)$. The GR identity $d\tau = \sqrt{g_{00}}\, dt$ is the standard-physics statement of the substrate-level becoming-rate ratio.

\textbf{Claim 7.} Clock networks are ED-geometry sensors. Differential measurements at $10^{-19}$ precision resolve millimeter-scale ED-gradient features on Earth.

\textbf{Claim 8 (negative).} No new substrate primitives are required beyond the seven listed in \S2.

\textbf{Claim 9 (scope-limit).} This walkthrough does not derive: many-body ED-flow beyond mean-field; nonlinear becoming gradients; strong-field relativistic corrections; quantum-coherence-enhanced clock precision; frequency-comb-based comparison; fundamental-constant drift; Lorentz/CPT-violation tests.

\textbf{The unified statement.} The Hau--Katori--Ye cluster is one substrate-level enterprise: manipulate or measure the local rate of becoming. The local rate of becoming is the rate of phase accumulation in pre-individuation amplitudes --- the substrate's foundational rate primitive. Speed of light, atomic-clock tick rate, refractive index, gravitational time dilation, and clock-network gradient sensing are all coarse-grained expressions of this same becoming-rate field. The substrate-level mechanism unifies AMO physics with general relativity at the level of the becoming-rate field.

\section*{References}

\begin{itemize}[leftmargin=*,itemsep=2pt]
  \item Hau, L. V., Harris, S. E., Dutton, Z., Behroozi, C. H. \emph{Light speed reduction to 17 metres per second in an ultracold atomic gas.} Nature \textbf{397}, 594 (1999).
  \item Liu, C., Dutton, Z., Behroozi, C. H., Hau, L. V. \emph{Observation of coherent optical information storage in an atomic medium using halted light pulses.} Nature \textbf{409}, 490 (2001).
  \item Katori, H. \emph{Spectroscopy of Strontium Atoms in the Lamb-Dicke Confinement.} Phys. Rev. Lett. \textbf{91}, 173005 (2003).
  \item Ye, J., Kimble, H. J., Katori, H. \emph{Quantum state engineering and precision metrology using state-insensitive light traps.} Science \textbf{320}, 1734 (2008).
  \item Bloom, B. J. et al. \emph{An optical lattice clock with accuracy and stability at the $10^{-18}$ level.} Nature \textbf{506}, 71 (2014).
  \item Bothwell, T. et al. \emph{Resolving the gravitational redshift across a millimetre-scale atomic sample.} Nature \textbf{602}, 420 (2022).
  \item Fleischhauer, M., Imamoglu, A., Marangos, J. P. \emph{Electromagnetically induced transparency: Optics in coherent media.} Rev. Mod. Phys. \textbf{77}, 633 (2005).
  \item Lukin, M. D. \emph{Trapping and manipulating photon states in atomic ensembles.} Rev. Mod. Phys. \textbf{75}, 457 (2003).
  \item Ludlow, A. D., Boyd, M. M., Ye, J., Peik, E., Schmidt, P. O. \emph{Optical atomic clocks.} Rev. Mod. Phys. \textbf{87}, 637 (2015).
  \item Misner, C. W., Thorne, K. S., Wheeler, J. A. \emph{Gravitation.} W. H. Freeman (1973).
\end{itemize}

\end{document}
